% =====================================================================
% paper/main-revised.tex  --  REVISED DRAFT FOR ADVISOR DISCUSSION
% Status: 2026-09-23. Structurally complete; experimentally INCOMPLETE.
%
% Rules this draft follows (see research/claim-audit.md section 5):
%  * source/main-original.tex is not modified (sha256 06f5f9d7...07cf6).
%  * No result for the proposed router is reported: none exists.
%  * The only measured numbers are the Stage-0 simple-baseline results in
%    Section VI. They are \input from paper/tables/*.tex, which
%    paper/make_tables.py generates from
%    experiments/prototype/results/stage0_raw.json (raw evidence:
%    experiments/prototype/{preds,splits,logs}).
%  * Every \cite key resolves to references/library.bib
%    (checked by paper/check_paper.py).
%  * Planned work is written in the future tense and marked \todo{...}.
%  * The title is a WORKING title. Whether to keep "Tiny-LM" and
%    "Jacobi-Forced Micro-Encoder" is an advisor decision
%    (research/micro-encoder-investigation.md, research/gap-assessment.md).
% Compile: cd paper && pdflatex main-revised && bibtex main-revised &&
%          pdflatex main-revised && pdflatex main-revised
% (not compile-tested: no TeX installation on the drafting host).
% =====================================================================
\documentclass[lettersize,journal]{IEEEtran}
\usepackage{amsmath,amsfonts,amssymb}
\usepackage{array}
\usepackage{textcomp}
\usepackage{stfloats}
\usepackage{url}
\usepackage{graphicx}
\usepackage{hyperref}
\usepackage{cite}
\usepackage{pifont}
\newcommand{\cmark}{\ding{51}}
\newcommand{\xmark}{\ding{55}}
\newcommand{\pmark}{$\triangle$}
\newcommand{\todo}[1]{\textbf{[TODO: #1]}}
\newcommand{\A}{\mathcal{A}}
\hyphenation{op-tical net-works semi-conduc-tor}

\begin{document}

\title{Query Routing between Edge SLMs, Split-LLMs, and Cloud LLMs:\\
Problem Formulation, Prior-Art Positioning, and a Pre-Registered Evaluation Protocol\\
{\normalsize (working title; revised draft, no results for the proposed router yet)}}

\author{Brendon Pedro, Sirapop Nuannimnoi, Panuwat Saetee, Hojjat Baghban,
Karin Honda, and Ching-Yao Huang%
\thanks{Revised draft prepared 2026-09-23 with AI assistance for advisor
discussion; see advisor/research-sprint-report.md. Not for submission.}}

\markboth{Revised draft for advisor discussion, 2026-09-23}{Pedro \MakeLowercase{\textit{et al.}}: Split-aware query routing (draft)}

\maketitle

\begin{abstract}
Collaborative inference between small language models (SLMs) on edge devices and
large language models (LLMs) in the cloud is an active research area, and
query-level routing between the two tiers has been studied extensively in both
datacenter and edge--cloud settings~\cite{li2026collaborative,moslem2026dynamic}.
This draft revisits a proposal for a lightweight edge-side router that chooses,
for each query, among an edge SLM, a split (partitioned) LLM, and a full cloud
LLM. After a review of prior work, we find that lightweight learned routers,
cost- and latency-aware edge--cloud routing, adaptation to new models without
retraining, and privacy-aware three-way edge/collaborative/cloud routing have
all been reported previously. In the literature we could read, we did not find
a query-level router whose action set includes a layer-split LLM path, with the
path's cost including activation transmission and split-depth-dependent
leakage. One likely-relevant 2026 journal paper could not be read and may close
this gap. We therefore (i) formulate split-aware routing with an explicit action
set, a measured system-state-dependent cost, and a leakage constraint; (ii)
position the formulation against the closest prior work; (iii) pre-register a
staged evaluation protocol with falsification criteria; and (iv) report a
zero-inference preliminary measurement that uses released per-query scores from
LLMRouterBench~\cite{li2026llmrouterbench} to set the bar a learned router must
clear. In that measurement, a TF-IDF plus logistic-regression router did not
outperform random routing on deferral AUC in any of eight leave-one-dataset-out
folds. No result for the proposed router is reported in this version.
\end{abstract}

\begin{IEEEkeywords}
LLM routing, edge--cloud collaborative inference, small language models,
split inference, selective prediction, calibration
\end{IEEEkeywords}

% ---------------------------------------------------------------------
\section{Introduction}
\label{sec:intro}
% Rewritten: the original first abstract sentence reproduced the first sentence
% of the abstract of li2026collaborative verbatim (research/gap-assessment.md s5).

\IEEEPARstart{L}{anguage-model} services are increasingly deployed across a
device--edge--cloud continuum rather than in a single location. Recent surveys
organise this space into collaboration between on-device SLMs and cloud LLMs,
including routing, cascading, speculative decoding, and partitioned
execution~\cite{li2026collaborative,chen2026network,moslem2026dynamic}. Two
motivations recur. First, per-token cost and latency of large models make them
wasteful for queries that a small model answers correctly~\cite{ding2024hybrid,ong2025routellm}.
Second, some inputs are sensitive, and sending them to a third-party cloud raises
privacy concerns~\cite{he2025emerged,zhang2025understanding}; even de-identified
clinical free text carries a residual re-identification risk~\cite{ford2025patient}.
\todo{If a regulatory claim (e.g.\ HIPAA) is kept, cite the primary regulatory
text; the original draft cited a company blog.}

Split (partitioned) inference is one response to the second concern. The first
layers of a model run on a trusted device and only intermediate activations
leave it~\cite{kang2017neurosurgeon,mudvari2024splitllm,lee2026browser}. It must
be distinguished from split \emph{learning}, where the partition is used during
training~\cite{vepakomma2018split,pham2025split,hukkeri2025comprehensive,alhindi2025balancing}.
Split inference does not by itself make inference private: transmitted
activations can be inverted. One preprint reports that an inversion attack
recovered about 59\% of tokens at a 2-layer split and about 35\% at an 8-layer
split~\cite{cunningham2026privacy}. Noise-based defences such as
Split-and-Denoise~\cite{mai2024split} trade utility for protection.

Routing between an edge SLM and a cloud LLM is well studied. Pre-generation
routers predict from the query whether the small model will suffice~\cite{ding2024hybrid,ong2025routellm,qiao2026consroute};
confidence-based cascades let the small model answer first and escalate when
uncertain~\cite{aggarwal2023automix,wu2025recursive,zhu2026dabo,gui2026pro}. PRISM routes
among cloud-only, edge-only, and a cloud--edge collaboration mode based on
privacy sensitivity~\cite{zhan2025prism}. Lightweight encoder or linear routers
are common~\cite{ong2025routellm,chen2024routerdc,garg2026hydra,dey2026compliancegate},
and several methods integrate new models without retraining the
router~\cite{wang2026icl,tran2025arch,shirkavand2025cost,garg2026hydra}. A large
benchmark further reports that many recent routers do not reliably outperform a
simple baseline~\cite{li2026llmrouterbench}.

In the sources we could read, we did not find a router whose action set contains
a \emph{split} LLM path alongside an edge SLM and a full cloud LLM. The question
this draft studies is therefore narrow: \emph{when an edge device can answer
locally, run part of a large model locally and send activations, or send the
query to the cloud, can a lightweight router learn from query content which path
to take, beyond what a fixed systems rule achieves?} Whether this question is
still open depends on a 2026 journal paper we could not access, Mix-or-Split~\cite{wang2026mix},
whose title suggests a per-request choice between collaborative and split
execution (Section~\ref{sec:related:open}).

\subsection*{What this draft contributes, and what it does not}
\begin{itemize}
\item \textbf{A formulation} of split-aware query routing
(Section~\ref{sec:problem}): an explicit action set
$\{\textsc{Edge},\textsc{Split}(d),\textsc{Cloud}\}$, a learned utility estimate
separated from a measured, system-state-dependent cost so that the trade-off
$\lambda$ can be changed at decision time (a property shared with prior
routers), and a split-depth-dependent leakage constraint.
\item \textbf{A positioning} against the closest prior work
(Section~\ref{sec:related}, Table~\ref{tab:comparison}), stating which parts of
the original proposal are already reported.
\item \textbf{A pre-registered, staged evaluation protocol}
(Section~\ref{sec:protocol}) with hypotheses H1--H3 and success and failure
criteria fixed before any router is trained.
\item \textbf{A preliminary, zero-inference measurement} of the simple-baseline
bar on released LLMRouterBench scores (Section~\ref{sec:stage0}). It covers
baselines only.
\end{itemize}
This draft does \emph{not} report a trained router, on-device latency, split-link
measurements, cost savings, or leakage measurements. These remain future work
(Section~\ref{sec:results}).

\subsection*{On the phrase ``Jacobi-Forced Micro-Encoder''}
The original title named a ``Jacobi-Forced Micro-Encoder''. As of 2026-09-23 we
found no established method with this name. Jacobi Forcing~\cite{hu2025fast} is
a training procedure that makes an autoregressive model produce correct tokens
under Jacobi-style parallel decoding. It builds on Jacobi decoding~\cite{song2021accelerating,santilli2023accelerating},
lookahead decoding~\cite{fu2024break}, and consistency LLMs~\cite{kou2024cllms}.
Its benefit is fewer sequential steps when generating long outputs with a
\emph{causal} model. The routing decision in Eq.~\eqref{eq:rule} is a single
forward pass followed by an $\arg\max$, so parallel decoding has nothing to
accelerate there, and an encoder-only model lacks the causal structure the
method assumes. Section~\ref{sec:design:jacobi} records the one role in which a
Jacobi-family method is technically relevant here: accelerating the edge or
split \emph{generator}. That role is orthogonal to routing. \todo{Advisor:
decide whether the title keeps ``Tiny-LM'' and ``Jacobi''. H2
(Section~\ref{sec:protocol}) tests the first.}

% ---------------------------------------------------------------------
\section{Related Work}
\label{sec:related}

\subsection{Query-level LLM routing}
Hybrid LLM trains a router that predicts the quality gap between a small and a
large model and routes under a tunable quality--cost
threshold~\cite{ding2024hybrid}. RouteLLM learns weak/strong routers from
preference data. Its best routers report cost-saving ratios of about
1.4--3.7$\times$ depending on the benchmark, at 87--95\% of the strong model's
performance~\cite{ong2025routellm}.\footnote{The original draft stated
2--3.6$\times$; the range above follows Table~6 of~\cite{ong2025routellm}
(3.66$\times$ on MT-Bench, 1.41$\times$ on MMLU, 1.49$\times$ on GSM8K).
RouteLLM's threshold is swept to meet call-rate targets rather than calibrated to
measured costs.} RouterDC~\cite{chen2024routerdc} and CSCR~\cite{shirkavand2025cost}
use contrastive query encoders. IRT-Router models query difficulty and model
ability with item response theory~\cite{song2025irt}. ICL-Router represents
models as in-context vectors so that new models can be added without retraining
the router~\cite{wang2026icl}. BEST-Route jointly selects a model and a number of
samples~\cite{ding2025best}. Arch-Router is a 1.5B-parameter \emph{generative}
router that maps queries to natural-language routing policies. Its reported
93.17\% accuracy and 28$\times$ latency advantage were measured on the authors'
own benchmark against proprietary LLMs prompted as routers, not against trained
routers~\cite{tran2025arch}. RouteLMT predicts marginal gain from the small
model's own prompt representations for translation~\cite{luo2026routelmt}, and
a related preprint routes from prefill activations of an open LLM~\cite{varshney2026llm}.
In production, HyDRA reports a CPU-served ModernBERT-based router~\cite{garg2026hydra}.
LLMRouterBench evaluates many routers on a common pool and reports that several
recent methods do not reliably outperform a simple baseline~\cite{li2026llmrouterbench}.
FrugalGPT is an early LLM cascade~\cite{chen2023frugalgpt}, and a recent survey organises routing and cascading methods~\cite{moslem2026dynamic}.

\subsection{Edge--cloud SLM/LLM collaboration}
Edge--cloud routers include ConsRoute (device--edge--cloud, prefill-state
probes)~\cite{qiao2026consroute}, a wireless quality--latency-aware
router~\cite{bao2025dynamic}, TMO (device--cloud reinforcement-learning offloading
in multi-turn conversations)~\cite{yuan2025device}, Lyra (online preference-aware
edge--cloud routing, IEEE INFOCOM 2026)~\cite{ma2026preference}, joint offloading
and SLM caching~\cite{xu2026joint}, and Pro-Router (a pre-scorer plus a
token-confidence verifier)~\cite{gui2026pro}. DABO calibrates an edge SLM's
confidence and offloads under an accuracy constraint~\cite{zhu2026dabo}.
LLMSwitchBench is a public local-versus-cloud routing benchmark for smart-home
queries~\cite{liu2025llmswitchbench}. Other work collaborates at the token level:
CITER routes tokens between models~\cite{zheng2025citer}, and distributed
speculative decoding drafts on the device and verifies in the
cloud~\cite{zheng2026fast,xu2026dip,zhang2026pipelined}.\footnote{The original
draft said TSLT~\cite{zheng2026fast} ``improves throughput and reduces latency''.
Its abstract claims reduced uplink communication while \emph{maintaining}
latency and quality. It is a speculative-decoding method, not sequence-level
routing.} Privacy-aware routing is represented by PRISM~\cite{zhan2025prism},
ComplianceGate~\cite{dey2026compliancegate}, and local privacy
gatekeepers~\cite{uzor2025guarding,langiu2026privacy}.

\subsection{Split inference for LLMs}
Split-point selection has been studied for DNNs~\cite{kang2017neurosurgeon} and
for LLMs under latency, memory, or network constraints~\cite{chen2024adaptive,sung2025memory,mudvari2024splitllm,younesi2025splitwise,kafetzis2025large,djuhera2025joint,feng2025distributed}.
Browser-based split inference reports that per-token round-trip time does not
depend on the split point~\cite{lee2026browser}. Work on the privacy of split
inference includes noise injection~\cite{mai2024split}, transport
optimisation~\cite{zhou2025poster}, encrypted activation transfer with an authenticated KV cache~\cite{li2026efficient},
and inversion measurements combined with lookahead decoding over a WAN~\cite{cunningham2026privacy}.
These systems choose \emph{where to split one model}. None that we read chooses
per query \emph{whether} to split, as opposed to answering locally or sending the
query to the cloud.

\subsection{Decision models, calibration, and selective prediction}
A router is a selective predictor: it abstains from the cheap model when
uncertain~\cite{geifman2017selective}, and its probability should be
calibrated~\cite{guo2017calibration}. A commercial ``System One'' decision-model
API (Jev) advertises typed decisions with confidence
estimates~\cite{almeida2026jev,typesafe2026systemone,typesafe2026confidencerouting}.
These are vendor claims that we have not verified. The service is cloud-hosted,
so it cannot act as a pre-offload router on the device.

\subsection{Positioning}
\label{sec:related:open}
Table~\ref{tab:comparison} summarises the closest work. It is condensed from
our full comparison matrix (\texttt{research/novelty-matrix.md}), in which a
difference counts only if it changes what must be built or measured. The
\emph{Proposed} row describes targets, not results.

\begin{table*}[t]
\centering
\caption{Closest prior work. \cmark: present; \pmark: partial; \xmark: absent
in the material we read; ?: not determined; --: not applicable. ``Abs.'' marks
rows read at abstract depth only. Venue status follows
\texttt{references/library.bib}; ``pp'' means preprint (a venue named only in
arXiv comments is not treated as confirmed).}
\label{tab:comparison}
\scriptsize
\setlength{\tabcolsep}{3pt}
\begin{tabular}{llcccccccc}
\hline
Work & Status & Edge SLM & Query-level & Lightweight & Latency- & Comm.- & Privacy- & Split-LLM & Calibrated \\
 & & tier & routing & router & aware & aware & aware & action & confidence \\
\hline
PRISM~\cite{zhan2025prism} & pp & \cmark & \cmark & \cmark & \cmark & \pmark & \cmark & \xmark$^a$ & \pmark \\
ConsRoute~\cite{qiao2026consroute} & pp & \cmark & \cmark & \cmark & \cmark & \pmark & \xmark & \xmark & \pmark \\
Bao et al.~\cite{bao2025dynamic} & pp (wksp.) & \cmark & \cmark & \cmark & \cmark & \cmark & \xmark & \xmark & \xmark \\
TMO~\cite{yuan2025device} (abs.) & pp & \cmark & \cmark & ? & \cmark & \pmark & \xmark & \xmark & \xmark \\
Lyra~\cite{ma2026preference} (abs.) & INFOCOM'26 & \cmark & \cmark & ? & \cmark & ? & ? & \xmark$^b$ & ? \\
DABO~\cite{zhu2026dabo} (abs.) & IEEE Access & \cmark & \cmark$^c$ & ? & \cmark & ? & ? & \xmark$^b$ & \cmark \\
Pro-Router~\cite{gui2026pro} & pp & \pmark & \cmark & \cmark & \cmark & \pmark & \xmark & \xmark & \cmark \\
Hybrid LLM~\cite{ding2024hybrid} (abs.) & pp & \pmark & \cmark & ? & \xmark & \xmark & \xmark & \xmark & \pmark \\
RouteLMT~\cite{luo2026routelmt} (abs.) & ACL'26 Ind. & \pmark & \cmark & \cmark & \xmark & \xmark & \xmark & \xmark & \pmark \\
RouteLLM~\cite{ong2025routellm} & ICLR'25 & \xmark & \cmark & \cmark & \xmark & \xmark & \xmark & \xmark & \pmark \\
HyDRA~\cite{garg2026hydra} & pp & \xmark & \cmark & \cmark & \cmark & \xmark & \xmark & \xmark & \pmark \\
ICL-Router~\cite{wang2026icl} (abs.) & AAAI'26 & \xmark & \cmark & \xmark & \xmark & \xmark & \xmark & \xmark & \pmark \\
IRT-Router~\cite{song2025irt} (abs.) & ACL'25 & \xmark & \cmark & ? & \xmark & \xmark & \xmark & \xmark & \xmark \\
Arch-Router~\cite{tran2025arch} & pp & \xmark & \cmark & \pmark & \pmark & \xmark & \xmark & \xmark & \xmark \\
Cunningham~\cite{cunningham2026privacy} (abs.) & pp & \pmark & \xmark & -- & \cmark & \cmark & \cmark & \cmark$^d$ & \xmark \\
Jev~\cite{almeida2026jev} (vendor) & product & \xmark & \cmark & \xmark & \cmark & \xmark & \xmark & \xmark & \cmark$^e$ \\
Mix-or-Split~\cite{wang2026mix} & TPDS'26 & ? & ? & ? & ? & ? & ? & ?$^f$ & ? \\
\hline
\textit{Proposed (this draft)} & not built & target & target & target & target & target & target & target & target \\
\hline
\end{tabular}

\vspace{2pt}
\footnotesize{$^a$ The collaboration mode perturbs the prompt with local
differential privacy and refines a cloud-returned sketch; no model is
partitioned. $^b$ No split path is mentioned in the abstract. $^c$ The SLM runs a
probe inference first (post-probe, not pre-generation). $^d$ A split system with
no router and no SLM path. $^e$ Vendor claim, not independently verified.
$^f$ Closed access; unread. Its title suggests it may occupy the Split-LLM column.}
\end{table*}

On this evidence, a lightweight or generative router, joint accuracy--cost--latency
routing between an edge SLM and a cloud LLM, adaptation without retraining, and
privacy-aware three-way routing are \emph{not} open problems. What we did not
find is the Split-LLM column combined with query-level routing. That combination
is a genuine research question only if the split path differs from the cloud path
in something that query content predicts, for example which inputs are too
sensitive to leave the device in any form other than activations. If it does
not, the split-versus-cloud choice may be a fixed systems rule rather than a
learned decision (Section~\ref{sec:problem:caveat}). \todo{Read Mix-or-Split
\cite{wang2026mix} in full (needs IEEE Xplore access) before relying on this
paragraph.}

% ---------------------------------------------------------------------
\section{Problem Formulation}
\label{sec:problem}

\begin{table}[t]
\centering
\caption{Notation}
\label{tab:notation}
\begin{tabular}{cp{6.2cm}}
\hline
Symbol & Meaning \\
\hline
$q$ & a user query \\
$s$ & observed system state (link bandwidth $B$, round-trip time $\rho$, cloud load) \\
$\A$ & action set $\{\textsc{Edge}_j\}_{j} \cup \{\textsc{Split}(d)\}_{d\in D} \cup \{\textsc{Cloud}\}$ \\
$D$ & admissible split depths (number of layers kept on the device) \\
$U(a,q)$ & quality of the answer produced by action $a$ (task metric, e.g.\ accuracy) \\
$C(a\mid q,s)$ & cost of action $a$ in one declared unit (Section~\ref{sec:problem:cost}) \\
$L(a)$ & leakage of action $a$ (fraction of input tokens recoverable by a stated attack) \\
$\varepsilon(q)$ & leakage budget for query $q$ \\
$\lambda$ & price of one cost unit in quality units \\
$\hat U_\theta(a\mid q)$ & router's estimate of $U(a,q)$, parameters $\theta$ \\
$\hat p_\theta(q)$ & router's probability that the edge answer is good enough \\
\hline
\end{tabular}
\end{table}

\subsection{Actions}
For each query the edge device chooses one action $a\in\A$: answer with an edge
SLM (\textsc{Edge}$_j$); run the first $d$ layers of a large model locally and
send activations to the cloud for the rest (\textsc{Split}$(d)$); or send the
query text to the cloud model (\textsc{Cloud}). The original draft named
Split-LLMs only as an option at the cloud tier and did not model the partition.
Here $d$ is an explicit decision variable.

\subsection{Objective}
The original objective, $\max_R \mathbb{E}_q[U(R(q),q)-\lambda C(R(q),q)]$,
follows the quality--cost trade-off used in prior routing
work~\cite{ding2024hybrid,ong2025routellm,song2025irt,ding2025best}. We keep it
and add system state and a leakage constraint:
\begin{equation}
\max_{\pi}\ \mathbb{E}_{q,s}\big[\,U(\pi(q,s),q) - \lambda\, C(\pi(q,s)\mid q,s)\,\big]
\ \ \text{s.t.}\ \ L(\pi(q,s)) \le \varepsilon(q).
\label{eq:objective}
\end{equation}
The router learns only the quantity that depends on query content. The cost is
measured, not learned:
\begin{equation}
\pi(q,s) = \arg\max_{a\in\A:\,L(a)\le\varepsilon(q)} \ \hat U_\theta(a\mid q) - \lambda\, C(a\mid q,s).
\label{eq:rule}
\end{equation}
Because $\lambda$, $s$, and the cost table enter only at decision time,
changing prices, links, or budgets does not require retraining. This is a
property of the decision rule, and prior routers share
it~\cite{ong2025routellm,ding2024hybrid}. It is not a contribution. Adding a new
\emph{model} does require new labels for $\hat U_\theta$ unless a method such
as~\cite{wang2026icl} is used.

In the two-action case (\textsc{Edge} vs.\ \textsc{Cloud}), Eq.~\eqref{eq:rule}
reduces to thresholding $\hat p_\theta(q)=\Pr[U(\textsc{Edge},q)\ge U(\textsc{Cloud},q)-\delta\mid q]$
at a threshold $\tau(\lambda,s)$. How $\tau$ is set on a new query distribution
must be specified: Section~\ref{sec:stage0} shows that thresholds tuned on
development data can fail under dataset shift.

\subsection{Cost model}
\label{sec:problem:cost}
$C$ is expressed in one declared unit. We propose end-to-end latency in seconds,
with monetary cost reported separately using the cost models of
Section~\ref{sec:protocol:metrics}. For generating $N$ tokens:
\begin{align}
T(\textsc{Edge}_j) &= T^{\text{pre}}_{j}(q) + (N-1)\,T^{\text{tok}}_{j},\\
T(\textsc{Cloud}) &= \rho + |q|/B + T^{\text{pre}}_{c}(q) + (N-1)\,T^{\text{tok}}_{c},\\
T(\textsc{Split}(d)) &= T^{\text{pre}}_{e,d}(q) + \tfrac{|q|\,h\,b}{B} + \rho + T^{\text{pre}}_{c,S-d}(q) \nonumber\\
&\quad + (N-1)\big(T^{\text{tok}}_{e,d} + \rho + \tfrac{h\,b}{B} + T^{\text{tok}}_{c,S-d}\big),
\end{align}
where $h$ is the hidden size, $b$ the bytes per transmitted element, $S$ the
number of layers, and subscripts $e,c$ denote the device and cloud parts. The
per-token round trip in the last line is the reason Jacobi-family parallel
decoding has been applied to split inference~\cite{cunningham2026privacy}. This
is an analytic model. Every term is to be \emph{measured} on the testbed
(Stage~2), not assumed. \todo{Measure; the original testbed description is
incomplete (cloud server unspecified).}

\subsection{Leakage}
$L(\textsc{Cloud})$ is the whole query and $L(\textsc{Edge}_j)=0$. For
$\textsc{Split}(d)$, $L$ is defined operationally as the token-recovery rate of
a stated inversion attack at depth $d$. Published numbers are
depth-dependent~\cite{cunningham2026privacy}. $\varepsilon(q)$ requires a
sensitivity label per query, as in PRISM~\cite{zhan2025prism}. \todo{Choose the
attack and the sensitivity labelling procedure.}

\subsection{When is the split action a learning problem?}
\label{sec:problem:caveat}
If $\textsc{Split}(d)$ runs the same weights as $\textsc{Cloud}$ and the
partition is lossless, then $U(\textsc{Split}(d),q)=U(\textsc{Cloud},q)$ for
every $q$. Quality labels alone then cannot distinguish the two actions; only
$C$ and $L$ can. In that case the split-versus-cloud choice is determined by
$s$ and $\varepsilon(q)$, and the router only needs to predict $\varepsilon(q)$
(sensitivity) and $\hat U(\textsc{Edge}_j\mid q)$. The split action requires
learning beyond this only if (a) the split path uses a different model, depth,
or lossy compression, so that quality differs, or (b) sensitivity must itself be
inferred from content. Hypothesis H3 (Section~\ref{sec:protocol}) tests whether
a learned policy beats the fixed rule ``sensitive $\rightarrow$ best admissible
\textsc{Split}$(d)$, otherwise cost-optimal''.

% ---------------------------------------------------------------------
\section{Candidate Router Designs}
\label{sec:design}
The original draft specified a ``Tiny Language Model'' router but left the method
section empty. We compare four candidate designs. None has been implemented or
evaluated yet. Details are in \texttt{architecture/options.md}.

\begin{enumerate}
\item[A.] \textbf{Tiny generative router.} A small instruction-tuned LM
(e.g.\ 0.5B parameters) fine-tuned to emit a route label; $\hat p$ is read from
the label-token probabilities. This is the closest reading of the original
title. Arch-Router~\cite{tran2025arch} is a larger instance.
\item[B.] \textbf{Discriminative encoder router.} A small encoder (e.g.\ a
DeBERTa-v3-small or ModernBERT-base backbone) with a classification head. This is the
design of RouteLLM's BERT router~\cite{ong2025routellm}, RouterDC~\cite{chen2024routerdc},
and HyDRA~\cite{garg2026hydra}.
\item[C.] \textbf{Self-routing edge SLM.} A probe on the edge SLM's own prefill
state, as in RouteLMT~\cite{luo2026routelmt}, ConsRoute~\cite{qiao2026consroute},
and~\cite{varshney2026llm}. It reuses computation the edge must do anyway
if it answers locally.
\item[D.] \textbf{Structured decision head.} Designs A--C with a typed output
(action, split depth, confidence, abstain), in the spirit of typed decision
APIs~\cite{typesafe2026systemone}. It runs on the device.
\end{enumerate}
Eq.~\eqref{eq:rule} needs one forward pass and an $\arg\max$. No design
requires sequential decoding. Design A must therefore be justified empirically
against B (H2).

\subsection{Where a Jacobi-family method could apply}
\label{sec:design:jacobi}
Jacobi Forcing trains a causal LM so that Jacobi (fixed-point) parallel decoding
converges in fewer iterations. Its authors report speedups for 7B code and math
models on datacenter GPUs, and the released checkpoints are two 7B
models~\cite{hu2025fast,jacobiforcing2025repo}. In this system it could
plausibly reduce the number of sequential edge--cloud round trips when the
\emph{split generator} decodes, the $(N-1)$ term in $T(\textsc{Split}(d))$,
analogously to the use of lookahead decoding in~\cite{cunningham2026privacy}.
Whether it transfers to sub-2B edge models or to a split link has not been
reported, to our knowledge. That would be a decoding contribution separate
from routing. \todo{Advisor: in scope or not?}

% ---------------------------------------------------------------------
\section{Evaluation Protocol (Proposed)}
\label{sec:protocol}
This section describes experiments that \emph{have not been run}, except
Stage~0's simple baselines (Section~\ref{sec:stage0}). The full pre-registration
is in \texttt{architecture/recommended-experiment.md}.

\subsection{Hypotheses}
\begin{itemize}
\item \textbf{H1 (a learned router beats simple baselines).} On held-out queries,
a learned pre-generation router reaches a lower cloud-call rate at quality
$\ge$ strong $-1$ point than both TF-IDF+LR and random routing at a matched
rate.
\item \textbf{H2 (a tiny generative router is justified).} Design A is not
significantly worse than design B and offers something B does not (e.g.\ better
out-of-distribution calibration). Otherwise the ``Tiny-LM'' framing is not
supported.
\item \textbf{H3 (the split action matters).} A learned policy over
$\{\textsc{Edge},\textsc{Split}(d),\textsc{Cloud}\}$ beats the fixed rule of
Section~\ref{sec:problem:caveat} on the objective in Eq.~\eqref{eq:objective}.
\end{itemize}

\subsection{Stages}
\textbf{Stage 0 (no model inference).} Labels come from the per-query scores
released with LLMRouterBench~\cite{li2026llmrouterbench,llmrouterbench2025repo}.
The pool is the binary-scored splits for which both models in a pair have
records (MMLU-Pro, GPQA, LiveCodeBench, LiveMathBench, AIME; 2,432 unique
queries). There are two pairs: P-L (Llama-3.1-8B-Instruct $\rightarrow$
Qwen3-235B-A22B-2507) and P-Q (Qwen3-8B $\rightarrow$ Qwen3-235B-A22B-2507).
Splits are 60/20/20 stratified with five seeds, plus four leave-one-dataset-out
(LODO) folds. \emph{Status: simple baselines run; designs A--C not run.}

\textbf{Stage 1 (own labels).} We will run the edge models of the original
draft on the Stage-0 datasets and on SimpleQA~\cite{wei2024measuring} to obtain
labels. The original draft assumed that LLMRouterBench contains SimpleQA scores
for its model pool. The released data contain SimpleQA scores only for
flagship API models, and of the original pool only Llama-3.1-8B-Instruct appears
in the release. Flagship models score low on SimpleQA in the release, so it offers
little headroom for SLM-versus-LLM routing, and a second dataset such as
LLMSwitchBench~\cite{liu2025llmswitchbench} should be considered. \emph{Status:
not run.}

\textbf{Stage 2 (split testbed).} We will measure the cost terms of
Section~\ref{sec:problem:cost} and the leakage per depth on the edge server and
cloud server, and test H3. \emph{Status: not run; cloud hardware unspecified.}

\subsection{Models}
The original model list needs corrections (verified against Hugging Face model
records on 2026-09-23; see \texttt{experiments/reproducibility-plan.md}).
No official ``Qwen3-72B-Instruct'' checkpoint exists. The official Gemma-4
identifiers are \texttt{gemma-4-E2B-it}, \texttt{gemma-4-E4B-it}, and
\texttt{gemma-4-12B-it}, not ``E2B-A4B'' etc. The E2B checkpoint stores about
5.1B parameters. The official Nemotron name includes ``A3B''. Candidate edge models
are Gemma-4 E2B/E4B/12B, Phi-4-mini~\cite{microsoft2025phi},
Llama-3.1-8B-Instruct, and the diffusion LMs LLaDA~\cite{nie2025large}, Dream,
and Fast-dLLM v2~\cite{wu2025fast}. Candidate strong models are
Llama-3.3-70B-Instruct and a replacement for the non-existent Qwen3-72B.
\todo{Advisor: choose the strong-model replacement (e.g.\ Qwen2.5-72B-Instruct,
Qwen3-32B, or Qwen3-235B-A22B-2507, which matches Stage 0).}

\subsection{Baselines}
Always-edge, always-cloud, random at a matched call rate, a dataset-prior router
(which routes on dataset identity and serves as a leakage control), TF-IDF+LR,
RouteLLM's released BERT router applied without retraining~\cite{ong2025routellm,routellm2024repo},
Arch-Router applied zero-shot~\cite{tran2025arch}, a prefill probe (design C),
and an oracle. RouteLMT's router code is not public~\cite{luo2026routelmt}, so it
will be re-implemented from the paper and described as such. ICL-Router and
IRT-Router~\cite{wang2026icl,song2025irt} belong to the full study. For H3: the
fixed split rule of Section~\ref{sec:problem:caveat}, and PRISM's
gate~\cite{zhan2025prism}.

\subsection{Metrics}
\label{sec:protocol:metrics}
\emph{Routing:} cloud-call rate at quality $\ge$ strong $-1$ point (primary),
deferral AUC (area under quality vs.\ call rate), ECE, Brier score, and AURC.
ECE is always reported next to a ranking metric, because a router can be
well calibrated and still uninformative (Section~\ref{sec:stage0}).
\emph{System:} router latency (p50/p95) on the device; TTFT, TPOT, and
$T_{\text{total}}=\text{TTFT}+(N-1)\,\text{TPOT}$; activation bytes per token;
requests and tokens per second; cold-start latency for serverless deployment.
\emph{Cost:} as in the original draft, dedicated-GPU cost per request
$C_{\text{cloud}}=P_{\text{gpu}}T_{\text{alloc}}/N_{\text{req}}$, serverless cost
$C_{\text{serverless}}=P_{\text{flex}}T_{\text{exec}}+C_{\text{overhead}}$, and edge
cost $C_{\text{edge}}=E_{\text{query}}R_{\text{elec}}+C_{\text{device}}/L$. Every
reported quantity will be labelled as measured or estimated (e.g.\ prices taken
from a price list). \todo{Record price sources and access dates.}

\subsection{Statistics and decision criteria}
The protocol specifies a paired bootstrap (10,000 resamples, stratified by
dataset), McNemar's exact test at the selected operating point,
Holm--Bonferroni correction across routers, pairs, and folds, and five seeds.
H1 is supported only if the best learned router lowers the primary metric by at
least 5 points relative to TF-IDF+LR (95\% CI excluding 0) and beats random and
dataset-prior on deferral AUC in at least 3 of 4 LODO folds. The thresholds are
our proposals and await the advisor's confirmation.

% ---------------------------------------------------------------------
\section{Preliminary Measurement: the Simple-Baseline Bar (Stage 0)}
\label{sec:stage0}
This section reports the only measurements in this draft. They were computed
from released per-query scores. No language model was run, and \emph{no proposed
router was trained}. Code, frozen splits, per-query predictions, logs, and
checksums are in \texttt{experiments/prototype/}. Tables~\ref{tab:stage0}
and~\ref{tab:stage0lodo} are generated from the raw results file by a script.

\begin{table}[t]
\centering
\caption{Stage 0, in-domain test (mean$\pm$std over 5 seeds; $n\approx486$ per
seed). ``Target met'': seeds where test quality reached strong $-1$ point at the
threshold tuned on the development split. $^\dagger$Routes on dataset identity
(a leakage control). $^\ddagger$Upper bound. Measured on released
LLMRouterBench scores; no model inference.}
\label{tab:stage0}
\scriptsize
\setlength{\tabcolsep}{3pt}
% AUTO-GENERATED by paper/make_tables.py from stage0_raw.json (sha256 3f83fff0260f7896069c363e5c966fd5864cca01fc25902f99c1de465dc24793). Do not edit.
\begin{tabular}{lccccc}
\hline
Router & Call rate & Quality & Target met & Deferral AUC & ECE \\
\hline
\multicolumn{6}{l}{\textit{Pair P-L}} \\
Always-edge & 0.000$\pm$0.000 & 0.298$\pm$0.000 & -- & -- & -- \\
Always-cloud & 1.000$\pm$0.000 & 0.719$\pm$0.000 & -- & -- & -- \\
Random & 0.972$\pm$0.013 & 0.709$\pm$0.004 & 2/5 & 0.509$\pm$0.007 & 0.258$\pm$0.017 \\
Dataset-prior$^\dagger$ & 1.000$\pm$0.000 & 0.719$\pm$0.000 & 5/5 & 0.529$\pm$0.001 & 0.002$\pm$0.000 \\
TF-IDF+LR & 0.949$\pm$0.014 & 0.709$\pm$0.006 & 3/5 & 0.542$\pm$0.003 & 0.052$\pm$0.010 \\
Oracle$^\ddagger$ & 0.441$\pm$0.000 & 0.739$\pm$0.000 & 5/5 & 0.636$\pm$0.001 & 0.000$\pm$0.000 \\
\multicolumn{6}{l}{\textit{Pair P-Q}} \\
Always-edge & 0.000$\pm$0.000 & 0.680$\pm$0.000 & -- & -- & -- \\
Always-cloud & 1.000$\pm$0.000 & 0.718$\pm$0.000 & -- & -- & -- \\
Random & 0.580$\pm$0.207 & 0.700$\pm$0.014 & 1/5 & 0.698$\pm$0.008 & 0.396$\pm$0.011 \\
Dataset-prior$^\dagger$ & 0.410$\pm$0.000 & 0.732$\pm$0.000 & 5/5 & 0.716$\pm$0.001 & 0.002$\pm$0.000 \\
TF-IDF+LR & 0.299$\pm$0.086 & 0.722$\pm$0.009 & 5/5 & 0.723$\pm$0.005 & 0.017$\pm$0.011 \\
Oracle$^\ddagger$ & 0.113$\pm$0.000 & 0.794$\pm$0.000 & 5/5 & 0.750$\pm$0.003 & 0.000$\pm$0.000 \\
\hline
\end{tabular}

\end{table}

\begin{table}[t]
\centering
\caption{Stage 0, deferral AUC on the held-out dataset group (LODO, seed 0).}
\label{tab:stage0lodo}
\scriptsize
% AUTO-GENERATED by paper/make_tables.py from stage0_raw.json (sha256 3f83fff0260f7896069c363e5c966fd5864cca01fc25902f99c1de465dc24793). Do not edit.
\begin{tabular}{llcccc}
\hline
Pair & Held-out & Random & Dataset-prior & TF-IDF+LR & Oracle \\
\hline
P-L & MMLU-Pro & 0.645 & 0.640 & 0.644 & 0.765 \\
P-L & GPQA & 0.427 & 0.418 & 0.430 & 0.560 \\
P-L & LiveCodeBench & 0.409 & 0.402 & 0.392 & 0.530 \\
P-L & math & 0.468 & 0.455 & 0.445 & 0.579 \\
P-Q & MMLU-Pro & 0.765 & 0.754 & 0.751 & 0.824 \\
P-Q & GPQA & 0.577 & 0.582 & 0.568 & 0.656 \\
P-Q & LiveCodeBench & 0.659 & 0.658 & 0.667 & 0.699 \\
P-Q & math & 0.769 & 0.752 & 0.746 & 0.789 \\
\hline
\end{tabular}

\end{table}

\textbf{Headroom depends on the pair.} With Llama-3.1-8B-Instruct as the edge
model (P-L), always-edge scored 0.298 against 0.719 for always-cloud. The oracle
met the target while calling the cloud for 44.1\% of queries. With Qwen3-8B
(P-Q), always-edge scored 0.680 against 0.718, and the oracle called the cloud
for 11.3\% of queries.

\textbf{The bar a learned router must clear.} In-domain, TF-IDF+LR met the
target in 3/5 seeds for P-L at a call rate of 0.949$\pm$0.014, and in 5/5 seeds
for P-Q at 0.299$\pm$0.086. Its deferral AUC was significantly higher than
random in every P-L seed and in 4 of 5 P-Q seeds. It was significantly higher
than the dataset-prior control in only 2 of 5 seeds for each pair (paired
bootstrap, 95\% CI).

\textbf{Lexical routing did not transfer across datasets.} In the LODO folds,
TF-IDF+LR was never significantly better than random on deferral AUC. It was
significantly worse in 3 of 8 folds (P-L LiveCodeBench, P-Q MMLU-Pro, P-Q math).

\textbf{Calibration alone is misleading.} The dataset-prior router had in-domain
ECE of 0.002 but little ranking ability (Table~\ref{tab:stage0}).

\textbf{Development-tuned thresholds can fail under shift.} In the P-Q MMLU-Pro
LODO fold, always-edge already met the target on the development groups, so the
tuned threshold sent nothing to the cloud and the target was missed on test,
even for the oracle's scores.

\textbf{Caveats.} These are two model pairs on five benchmarks, with scores
produced by the benchmark authors. Qwen3-8B's long outputs suggest a reasoning
mode, which we inferred but did not confirm. The released cost field covers API
cost only. The results say nothing about any on-device router, about latency, or
about the split action.

% ---------------------------------------------------------------------
\section{Results and Discussion}
\label{sec:results}
\todo{No results for the proposed router exist yet. This section will report
H1 and H2 (Stage 0, designs A--C), Stage-1 routing results on the original
edge-model pool, and Stage-2 latency, cost, communication, and leakage
measurements with H3. It must not be filled before the corresponding raw outputs
exist.}

\subsection{Limitations of this draft}
\begin{itemize}
\item The literature review was bounded. IEEE Xplore and the ACM Digital Library
could not be queried directly. Several rows of Table~\ref{tab:comparison} were
read at abstract depth only, and Mix-or-Split~\cite{wang2026mix} is unread.
\item The only measurements are simple baselines on released scores for two
model pairs, which do not include the original draft's edge models apart from
Llama-3.1-8B-Instruct.
\item The Split-LLM question may reduce to a fixed systems rule
(Section~\ref{sec:problem:caveat}).
\item The leakage measure and sensitivity labels are not yet defined.
\end{itemize}

% ---------------------------------------------------------------------
\section{Conclusion}
\label{sec:conclusion}
This revision narrows the original proposal. Lightweight, adaptive, cost- and
latency-aware SLM/LLM routing is established in prior work and is not claimed
here. What we did not find in the literature we could read is a query-level router whose
action set includes a split LLM path with measured activation cost and
depth-dependent leakage. Whether that question is open, and whether it requires
a learned router at all, remains to be established. We give a formulation that
makes it testable, a pre-registered protocol, and a measured simple-baseline
bar. Any claim of benefit awaits the experiments in Section~\ref{sec:protocol}.

\bibliographystyle{IEEEtran}
\bibliography{../references/library}

\end{document}
